\section{Data}
\label{sec:data}

\begin{table}[t]
\centering
\small
\setlength{\tabcolsep}{2.4pt}
\begin{tabular}{lrrrrr}
\toprule
Set & Utts & Hours & CS\% & \Emb{} units & Spk \\
\midrule
ASCEND train & 9{,}869 & 8.8 & 28 & 7{,}081 & 18 \\
ASCEND val & 1{,}130 & 0.9 & 18 & 471 & 3 \\
ASCEND test & 1{,}315 & 0.9 & 28 & 772 & 2 \\
NSC train & 35{,}647 & 40.0 & 62 & 105{,}493 & 412 \\
NSC val & 3{,}367 & 3.0 & 25 & 4{,}099 & 60 \\
NSC test (Same) & 8{,}876 & 8.0 & 29 & 11{,}042 & 120 \\
NSC test (Diff) & 4{,}085 & 4.0 & 29 & 5{,}648 & 626 \\
Ours (phone) & 120 & 0.18 & 100 & 227 & 3 \\
\bottomrule
\end{tabular}
\caption{Data. CS\% = share of utterances that contain a switch; Spk = speakers. NSC = NSC Part~4, Mandarin speakers; Same = Same Room, Diff = Different Room (Zoom). Ours = our phone recordings (test only).}
\label{tab:data}
\end{table}

\paragraph{Corpora.}
ASCEND~\citep{lovenia2022ascend} is spontaneous Mandarin--English conversation; 69\% of its sentences are Mandarin-main. NSC Part~4~\citep{koh2019nsc} is conversational Singapore speech in which speakers were asked to switch. We keep its Mandarin-speaking participants, found from the speaker information sheet; 70\% of their sentences are English-main, but about 4 in 10 of their code-switched sentences are Mandarin-main, so NSC also contains many embedded English words. NSC has two recording conditions: \emph{Same Room}, both speakers in one room, and \emph{Different Room}, over Zoom. Table~\ref{tab:data} gives the splits.

\paragraph{Sampling and splits.}
NSC training is capped at 40\,h of the 193\,h of usable Same Room speech: code-switched utterances first (up to 30\,h), then English-only ones. Validation and test are random samples within a session-level speaker split. The Different Room test is a random 4\,h sample, and none of its speakers appears in training. ASCEND validation and test speakers do not appear in training either. NSC code-switched utterances have more switches than ASCEND ones (2.85 vs.\ 2.12 per utterance, Different Room) and longer embedded spans (2.8 vs.\ 2.0 units).

\paragraph{Our phone recordings.}
To test critical words in real voices, we wrote 120 code-switched sentences dense in names, numbers, drug and course terms (\emph{my \zh{密码} has twelve characters now}), 60 English-main and 60 Mandarin-main. Three authors, English-speaking Singaporeans of Chinese background, each read a different 40 (20 of each kind), once, on their own phone in a quiet room. The 48\,kHz files go through the same pipeline as the corpora (16\,kHz mono); clips last 3.8--7.2\,s, 10.7\,min in total. Every clip follows its sentence exactly, so the sentence is the reference.

\paragraph{Labels.}
Word times come from forced alignment with MMS~\citep{pratap2024mms} on romanised text~\citep{hermjakob2018uroman}; on FLEURS~\citep{conneau2023fleurs} clips spliced at known points, the aligned switch is a median of 10\,ms from the true one. An embedded word replaced by text in the other script counts as a \emph{translation} if the pair is in CC-CEDICT~\citep{cccedict} or its LaBSE similarity~\citep{feng2022labse} is at least $\tau=0.65$. Otherwise it counts as a \emph{transliteration} (a sound-alike) if its phonetic similarity (pinyin and Metaphone codes compared by edit distance) is at least $\theta=0.1$. Both thresholds were chosen on 1{,}768 ASCEND pairs labelled by two LLM annotators and adjudicated ($\kappa=0.84$ for translation vs.\ the rest). The critical-word list combines a rule tagger with a list labelled the same way ($\kappa=0.77$). Two small human audits of these labels are reported in Limitations.
